\chapter{How to Run the Code}

Every algorithm in this book ships as a small, \emph{self-contained} Python
listing: each one embeds all of its helper routines (MST, matching, max-flow,
LP solver, tree traversal, etc.), so none of them require any external package
beyond the Python standard library.  A matching, fully verified implementation
of every chapter lives under \texttt{src/pycodes/}.

\section*{Prerequisites}

\begin{itemize}
  \item Python 3.8 or later (no third-party packages needed).
  \item \texttt{src/pycodes/} on the Python path if you want the full,
        tested modules rather than the listing text.
\end{itemize}

\section{Quick start}

Run \emph{all} chapter demos at once:

\begin{lstlisting}[style=python, caption=Run every verified demo]
$ cd src/pycodes
$ python3 main.py
ALL DEMOS COMPLETED
\end{lstlisting}

Run a single chapter's demo directly:

\begin{lstlisting}[style=python, caption=Run one chapter's demo]
$ python3 src/pycodes/set_cover.py
$ python3 src/pycodes/steiner_tsp.py
$ python3 src/pycodes/multiway_kcut.py
\end{lstlisting}

\section{Trying a book listing}

The listings in the chapters are written to be copied and pasted into a
research editor as-is.  Because they are self-contained, you can wrap the body
of any listing with a small driver and run it directly:

\begin{lstlisting}[style=python, caption=Run a standalone chapter listing]
# paste the full listing body from the chapter here ...
graph = {0:{1:1,2:3}, 1:{0:1,2:2}, 2:{0:3,1:2}}
for k in (1,2):
    centers, radius = kcenter_parametric_pruning(graph, k)
    print(k, centers, radius)
\end{lstlisting}

\begin{itemize}
  \item The book's listings are validated during the build: each one compiles
        without syntax errors, and the primary function of every chapter is
        exercised with a small example.
  \item If a listing references a helper (\texttt{mst\_prim}, \texttt{Simplex},
        \texttt{MaxFlow}, \dots), \emph{it is defined in the same listing}, so
        cutting and pasting the whole block always works.
\end{itemize}

\section{Directory layout}

\begin{lstlisting}[style=basic, caption=Where the code lives]
approx_algo.tex        # LaTeX source (master)
chapters/*.tex         # one file per chapter, 1-30
src/pycodes/*.py       # verified, runnable chapter implementations
src/cppcodes/          # equivalent C++23 library + examples (CMake)
\end{lstlisting}

\section{Rebuilding the book (optional)}

The book is plain \LaTeX{}.  Any standard engine works:

\begin{lstlisting}[style=basic, caption=Compile the book]
$ latexmk -pdf approx_algo.tex
\end{lstlisting}